\chapter*{Definizioni}
\addcontentsline{toc}{chapter}{Definizioni}

\noindent Questa tesi usa un piccolo vocabolario di termini greci, sanscriti e
indoeuropei ricostruiti in sensi precisi. Sono raccolti qui affinch\'e
l'argomentazione possa procedere senza interrompersi per glossarli.

\begin{description}

\item[\textsc{Ayu}] (sanscrito \emph{āyu}; protoindoeuropeo \emph{*h\textsubscript{2}óyu-},
``forza vitale, durata di vita assegnata, eternit\`a''; cfr.\ latino \emph{aevum},
inglese \emph{ever}, e \emph{eone}). La forza vivente che ci \`e data da spendere:
non la vita come fatto biologico ma la vita come dotazione finita di vitalit\`a,
un fondo d'azione. Il termine porta a un tempo \emph{vita} ed \emph{eternit\`a}
--- l'intuizione che una vita spesa nella coerenza raggiunga, attraverso ci\`o che
lascia dietro di s\'e, l'immortale. La tesi assume che \physis\ sia il
governatore di come questo \textsc{ayu} sia rettamente speso.

\item[\textsc{Theorein}] (greco \emph{θεωρεῖν}, ``contemplare, vedere''; radice di
\emph{teoria}). La tesi usa la parola nel suo senso \emph{pre-platonico}, recuperato
prima che Platone e Aristotele la volgessero nella serena speculazione delle forme
eterne. L'arcaico \emph{theoros} era un delegato sacro mandato all'oracolo o alla
festa, uno spettatore agli \emph{agones}, un viaggiatore (Solone parte ``per amore di
\emph{theoria}''): la parola nominava un \emph{servizio sacro}, una \emph{gara} e
un'\emph{esplorazione} fatti col corpo e sulla strada --- un uscire per testimoniare,
non uno sguardo da poltrona. Cos\`i qui \emph{theorein} non \`e lo spettare distaccato
ma il vedere che compie il conoscere attraverso il fare: il segmento pi\`u alto della
linea divisa (Libro VI) letto come atto, non come fissare. La revisione del proprio
grafo di coerenza \`e un \emph{theorein}, e l'intera prova vissuta (Poscritto) lo \`e
nel senso arcaico.

\item[\textsc{Aner -- nar -- h\textsubscript{2}ner -- vir}] la stirpe indoeuropea
dell'``uomo come portatore di forza vitale ed eroica.'' Greco \emph{anēr}
(\emph{ἀνήρ}, ``uomo, eroe,'' da cui \emph{andreia}, coraggio); sanscrito
\emph{nár-, nr̥} (uomo, potenza vitale); protoindoeuropeo
\emph{*h\textsubscript{2}nér-} (la forza maschile sovrana, il potere-di-vita che
\emph{agisce}); latino \emph{vir} (uomo, da cui \emph{virtus}, virt\`u). Il filo
nomina non un sesso ma una postura: l'\emph{agathos}, l'eccellente la cui natura
\`e agire bene e spendere il proprio \textsc{ayu} in atto cercante-forma. Il
praticante \`e l'erede di questa stirpe --- \emph{vir} divenuto \emph{virtus},
forza vitale divenuta eccellenza.

\item[\textsc{Agathos} / \textsc{Kleos}] \emph{agathos} (greco, ``buono, nobile,
eccellente''): colui la cui azione \`e eccellente, nel senso omerico e
aristocratico prima di quello moralizzato. \emph{kleos} (greco, ``gloria, ci\`o
che si ode''): la fama immortale che l'eroe conquista quando la sua impresa \`e
ricordata e ri-cantata. La tesi sostiene (Libro X) che l'UberWiki \`e lo strumento
moderno del \emph{kleos} --- gloria immortale tra i mortali --- perch\'e l'atto
coerente persiste come nodo che altri mettono in atto.

\item[\physis\ / \praxis\ / \techne] natura (il principio attivo di crescita e
auto-organizzazione) / azione (compiuta per se stessa, che esprime il carattere
dell'agente) / conoscenza-artigianale (l'arte del fare). La tesi: \praxis\ \`e lo
strumento di \physis, e \techne\ \`e giusta quando li serve, ingiusta quando li
usurpa.

\item[\textsc{Entelecheia}] (Aristotele) ``essere-all'opera-restando-s\'e'': lo
stato in cui una cosa ha il proprio fine dentro di s\'e e lo sta attivamente
realizzando. La forma che un essere \`e gi\`a all'opera per divenire.

\item[\textsc{Coerenza} ($+1$ / $0$ / $-1$)] il verdetto della natura su
un'azione: un effetto reale registrato ($+1$), nessun effetto ($0$), o
auto-violazione ($-1$). Non un punteggio morale ma una relazione misurata tra
intenzione dichiarata ed effetto registrato.

\item[\textsc{Olarchia} / \textsc{UberWiki}] l'unico grafo condiviso dell'azione
pesata per coerenza in cui ogni praticante deposita ci\`o che ha funzionato; il
bene comune del Libro X, l'\emph{anamnesi} delle macchine filosofiche.

\item[\emph{Das Man}] (Heidegger) il ``si'' anonimo --- lo standard pubblico
impersonale (``ci\`o che si fa'') che governa inautenticamente perch\'e nessuno
lo firma. La tesi tratta l'algoritmo di raccomandazione come \emph{das Man}
operazionalizzato, e l'UberWiki come la sua ri-autorizzazione.

\item[\textsc{Sogno}] il ciclo di sfondo del motore che rigioca i nodi inerti e
fallimentari e propone variazioni correttive --- la facolt\`a di variazione della
natura resa algoritmica.

\end{description}
